\documentclass[../../../include/open-logic-section]{subfiles}

\begin{document}

\olfileid{sth}{ordinals}{opps}
\olsection{ఉత్తరవర్తి, సీమా క్రమసంఖ్యలు}

తదుపరి కొన్ని అధ్యాయాల్లో
క్రమసంఖ్యలను విస్తృతంగా
వాడతాం. అందువల్ల కొన్ని
సరళమైన భావాలను
పరిచయం చేయడం ఉపయోగకరం.

\begin{defn}
ఏ క్రమసంఖ్య $\alpha$కైనా,
దాని \emph{ఉత్తరవర్తి}
$\ordsucc{\alpha}
=\alpha \cup \{\alpha\}$.
ఏదో క్రమసంఖ్య~$\beta$కు
$\ordsucc{\beta} = \alpha$
అయితే $\alpha$ను
\emph{ఉత్తరవర్తి}
క్రమసంఖ్య అంటాం.
$\alpha$ను \emph{సీమా}
క్రమసంఖ్య అంటాం
అప్పుడే మరియు అప్పుడే
$\alpha$ శూన్యం కాక,
ఉత్తరవర్తి క్రమసంఖ్యా
కాకపోతే.
\end{defn}
\noindent
ఇది \emph{ఉత్తరవర్తి}
అనే సరైన భావమని
ఈ ఫలితం చూపుతుంది:

\begin{prop}
ఏ క్రమసంఖ్య $\alpha$కైనా:
\begin{enumerate}
	\item $\alpha \in \ordsucc{\alpha}$;
	\item $\ordsucc{\alpha}$ ఒక క్రమసంఖ్య;
	\item $\alpha \in \beta \in
	\ordsucc{\alpha}$ అయ్యే క్రమసంఖ్య $\beta$ లేదు.
	\end{enumerate}
\end{prop}

\begin{proof}
స్పష్టంగా,
$\alpha \in \alpha \cup \{\alpha\} = \ordsucc{\alpha}$.
అలాగే $\ordsucc{\alpha}$
క్రమసంఖ్యల సంక్రామక
సమితి; కాబట్టి
\olref[basic]{corordtransitiveord}
ప్రకారం అది క్రమసంఖ్య.
$\alpha \in \beta \in \ordsucc{\alpha}$
అనేది అసాధ్యం;
ఎందుకంటే అప్పుడు
$\beta \in \alpha$
లేదా $\beta = \alpha$;
ఇది \olref[basic]{ordordered}కు
విరుద్ధం.
\end{proof}
\noindent
దీనితో అతిపరిమిత
ఆగమన నిరూపణకు
మరొక రూపం లభిస్తుంది:

\begin{thm}[సరళ అతిపరిమిత ఆగమనం]\ollabel{simpletransrecursion}
$\phi(x)$ ఈ షరతులను
పూర్తి చేసే సూత్రం
అనుకుందాం:
\begin{enumerate}
	\item $\phi(\emptyset)$; మరియు
	\item ఏ క్రమసంఖ్య $\alpha$కైనా, $\phi(\alpha)$ అయితే
	$\phi(\ordsucc{\alpha})$; మరియు
	\item $\alpha$ సీమా క్రమసంఖ్య అయి, $(\forall \beta \in
	\alpha)\phi(\beta)$ అయితే, $\phi(\alpha)$.
\end{enumerate}
అప్పుడు $\forall \alpha \phi(\alpha)$.
\end{thm}

\begin{proof}
విపర్యయాన్ని నిరూపిస్తాం.
కాబట్టి $\lnot\phi$
అయ్యే ఏదో క్రమసంఖ్య
ఉందనుకుందాం; అలాంటి
అత్యల్ప క్రమసంఖ్యను
$\gamma$ అనుకుందాం.
అప్పుడు $\gamma = \emptyset$
లేదా ఏదో $\alpha$కు
$\gamma = \ordsucc{\alpha}$,
అక్కడ $\phi(\alpha)$;
లేదా $\gamma$ సీమా
క్రమసంఖ్య, అలాగే
$(\forall
\beta \in \gamma)\phi(\beta)$.
ఈ మూడు సందర్భాలూ
సంబంధిత పరికల్పనలకు
విరుద్ధం.
\end{proof}
\noindent
చివరిగా, తరువాత
ఉపయోగపడే మరో సంకేతం:

\begin{defn}\ollabel{defsupstrict}
$X$ క్రమసంఖ్యల
సమితి అయితే,
$\supstrict(X) = \bigcup_{\alpha
\in X} \ordsucc{\alpha}$.
\end{defn}
\noindent
ఇక్కడ ``lsub'' అంటే
``కనిష్ఠ కఠిన పై
హద్దు''.\footnote{కొన్ని
పుస్తకాలు దీనికి
``$\text{sup}(X)$''
అనే సంకేతాన్ని వాడతాయి.
కానీ ఇతర పుస్తకాలు
కనిష్ఠ \emph{కఠినంకాని}
పై హద్దుకు, అంటే కేవలం
$\bigcup X$కు, అదే
``$\text{sup}(X)$''ను
వాడతాయి. $X$కు
గరిష్ఠ మూలకం $\alpha$
ఉంటే, ఈ రెండు
భావాలు విడిపోతాయి:
కనిష్ఠ \emph{కఠిన}
పై హద్దు
$\ordsucc{\alpha}$;
కనిష్ఠ \emph{కఠినంకాని}
పై హద్దు కేవలం
$\alpha$.} ఈ కింది
ఫలితం దానిని
వివరిస్తుంది:

\begin{prop}
$X$ క్రమసంఖ్యల
సమితి అయితే,
$\supstrict(X)$, $X$లోని
ప్రతి క్రమసంఖ్యకన్నా
ఎక్కువగా ఉండే
అత్యల్ప క్రమసంఖ్య.
\end{prop}

\begin{proof}
$Y = \Setabs{\ordsucc{\alpha}}{\alpha \in X}$
అనుకుందాం; అప్పుడు
$\supstrict(X) = \bigcup Y$.
క్రమసంఖ్యలు సంక్రామకాలు,
ప్రతి క్రమసంఖ్యలోని
మూలకమూ క్రమసంఖ్యే
కాబట్టి, $\supstrict(X)$
క్రమసంఖ్యల సంక్రామక
సమితి. అందువల్ల
\olref[basic]{corordtransitiveord}
ప్రకారం అది
క్రమసంఖ్య.

$\alpha \in X$ అయితే,
$\ordsucc{\alpha} \in Y$;
కాబట్టి $\ordsucc{\alpha}
\subseteq \bigcup Y = \supstrict(X)$;
అందువల్ల $\alpha \in
\supstrict(X)$. కాబట్టి
$\supstrict(X)$, $X$లోని
ప్రతి క్రమసంఖ్యకన్నా
కఠినంగా ఎక్కువ.

తిరిగి, $\alpha \in \supstrict(X)$
అయితే ఏదో $\beta \in X$కు
$\alpha \in
\ordsucc{\beta} \in Y$;
అందువల్ల $\alpha \leq
\beta \in X$. కాబట్టి
ఈ కుటుంబానికి ఉన్న ఏ కఠిన
పై హద్దూ ఈ సమితిలోని
ప్రతి మూలకంకన్నా
ఎక్కువగా ఉంటుంది;
$\supstrict(X)$, $X$
యొక్క \emph{కనిష్ఠ}
కఠిన పై హద్దు.
\end{proof}

\end{document}
